\section{From quotient classification to the Partition Principle}
\label{sec:conclusion}

We now pass from the classification of quotients of the label set to
arbitrary surjections.  The reduction uses only three statements inside
the symmetric model: every nonempty set is an image of the label set
times an ordinal, choice holds for ordinal-indexed families, and every
nonwellorderable image of the label set is bijective to it.  We first
prove this reduction in $\ZF$, independently of the construction.

The ordinal slicing and assembly of injections follow the local-to-global
method in Ryan-Smith's analysis of $\PP$
\cite[Propositions~3.17 and~3.20]{RyanSmith2025}. We use the stronger
classification of nonwellorderable images stated below and give the
complete reduction under these hypotheses.

\begin{lemma}[The all-set reduction]\label{lemma:pp-reduction}
Work in $\ZF$.  Suppose that a nonempty set $A$ has the following
properties.
\begin{enumerate}
\item For every nonempty set $X$, there are an ordinal $\eta$ and
a surjection $A\times\eta\longrightarrow X$.
\item Every ordinal-indexed family of nonempty sets has a choice
function.
\item Every nonwellorderable surjective image of $A$ is in bijection
with $A$.
\end{enumerate}
Then $\PP$ holds: for all sets $X,Y$, every surjection
$f:X\longrightarrow Y$ has an injection $j:Y\longrightarrow X$.
\end{lemma}

\begin{proof}
We will use two elementary observations about images of $A$.
If $h:A\longrightarrow B$ is onto and $\varnothing\ne S\subseteq B$,
then $S$ is also an image of $A$.  Indeed, fix one $s_0\in S$ and
define
\[
 h_S(a)=
 \begin{cases}
 h(a),&h(a)\in S,\\
 s_0,&h(a)\notin S.
 \end{cases}
\]
This proves an existence assertion for one given nonempty subset;
it does not choose a default element simultaneously for a family.
Also, the image of a wellorderable set under any function is
wellorderable: order its image by the least preimages in a fixed
wellordering of its domain.  Both observations are theorems of $\ZF$.

Fix a surjection $f:X\longrightarrow Y$.  If $Y=\varnothing$, the
empty function is the desired injection.  Otherwise $X$ is nonempty.
Choose a surjection $s:A\times\eta\longrightarrow X$ from the first
assumption.  Define, for $i<\eta$,
\begin{align}
 B_i&=s[A\times\{i\}],\notag\\
 Y_i&=f[B_i]\setminus\bigcup_{k<i}f[B_k],\notag\\
 S_i&=B_i\cap f^{-1}[Y_i].\label{eq:final-disjoint-slices}
\end{align}
Separation and Replacement form these sequences.  Every $y\in Y$
belongs to $f[B_i]$ for some $i$, and its least such index is the
unique index for which $y\in Y_i$.  Thus the $Y_i$ partition $Y$.
The $S_i$ are pairwise disjoint because their $f$-images lie in the
pairwise disjoint $Y_i$.  Moreover,
\[
 f\mathbin{\upharpoonright}S_i:S_i\longrightarrow Y_i
\]
is onto for every $i$.  The sets $S_i$ need not cover $X$.

Before choosing any injections, define
\begin{equation}\label{eq:final-candidate-injections}
 \begin{split}
 \mathcal J_i=\{j\in\mathcal P(Y\times X):{}&
    j\text{ is an injective function},\\
    &\dom(j)=Y_i,\quad\ran(j)\subseteq S_i\}.
 \end{split}
\end{equation}
Power Set and Separation give each candidate set, and Replacement
gives the sequence $\langle\mathcal J_i:i<\eta\rangle$.
We show that every candidate set is nonempty, fixing an arbitrary
$i<\eta$.

If $Y_i$ is empty, its empty injection belongs to $\mathcal J_i$.
Suppose next that $Y_i$ is nonempty and wellorderable.  Choose a
wellordering of this one set.  The nonempty fibers
\[
 S_i\cap f^{-1}[\{y\}]\qquad(y\in Y_i)
\]
can be indexed by an ordinal using that wellordering.  The second
assumption gives a choice function on them.  Since different fibers
are disjoint, this function is an injection $Y_i\longrightarrow S_i$.
It therefore belongs to $\mathcal J_i$.

Finally, suppose that $Y_i$ is nonwellorderable.  Then $S_i$ is
nonempty and nonwellorderable, by the second elementary observation
and the surjectivity of $f|S_i$.  The map $a\mapsto s(a,i)$
surjects from $A$ onto $B_i$.  The first observation makes the
nonempty subset $S_i\subseteq B_i$ an image of $A$, and composition
with $f|S_i$ makes $Y_i$ an image of $A$ as well.  By the third
assumption both $S_i$ and $Y_i$ are bijective to $A$.  Composing
one of these bijections with the inverse of the other gives an
injection $Y_i\longrightarrow S_i$, again in $\mathcal J_i$.

We have proved $(\forall i<\eta)\ \mathcal J_i\ne\varnothing$.
The pointwise proof selected neither a sequence of wellorderings nor
a sequence of quotient bijections.  Apply the second assumption now
to the already defined sequence of candidate sets.  It supplies
$\langle j_i:i<\eta\rangle$ with $j_i\in\mathcal J_i$.
Set
\[
 j=\bigcup_{i<\eta}j_i.
\]
The disjointness of the $Y_i$ makes this a function with domain $Y$.
If $j(y)=j(y')$, that value lies in both corresponding source pieces
$S_i$.  Their disjointness makes the indices equal, and injectivity
of that $j_i$ gives $y=y'$.  Thus $j:Y\longrightarrow X$ is an
injection.

Every construction in this proof takes place in the same model of
$\ZF$.  No rank bound on $X$ or $Y$, and no uniform choice of auxiliary
presentations for the slices, has been used.
\end{proof}

\begin{remark}\label{rem:quotients-versus-images}
The third hypothesis of Lemma~\ref{lemma:pp-reduction} can equivalently
be stated for quotients $A/E$ by equivalence relations on $A$.
For a surjection $h:A\longrightarrow B$, define
\[
 a\mathrel{E_h}a'\quad\Longleftrightarrow\quad h(a)=h(a').
\]
The map $[a]_{E_h}\mapsto h(a)$ is a bijection $A/E_h\longrightarrow B$.
Its graph is defined by the unique common value on each class; it
does not require choosing class representatives.

The first hypothesis is the surjective assertion $\mathsf{SVC}(A)$.
Its injective counterpart $\mathsf{SVC}^{+}(A)$ would assert an
injection $X\longrightarrow A\times\eta$ for every $X$ and some
ordinal $\eta$.  The proof does not assume that counterpart.
It follows afterward by applying $\PP$ to each presentation surjection.
The quotient assertion $\mathsf{PP}(A)$ says that each surjective
image of $A$ injects into $A$.  The subset-restricted assertion
$\PP\mathbin{\upharpoonright}A$ concerns surjections between two
arbitrary subsets of $A$.  These are different local statements.
The reduction above uses the stronger classification of every
nonwellorderable image of $A$, together with ordinal-indexed choice;
it does not replace that classification by quotient injections alone.
\end{remark}

The injection obtained in Lemma~\ref{lemma:pp-reduction} need not
respect the fibers of the original surjection.  In its wellorderable
target case the local injection is a section, but in its
nonwellorderable case the two bijections with $A$ give no such
requirement.  This is precisely the freedom allowed by $\PP$.

\subsection{Further cardinal-comparison consequences}

The same model separates Choice from two weaker cardinal-comparison
principles. Their consequences follow from $\PP$ alone and do not
require any further forcing construction.

\begin{remark}\label{rem:dual-csb-wpp}
For sets $X,Y$, write $X\leq^{*}Y$ if $X=\varnothing$ or there is a
surjection $Y\twoheadrightarrow X$, and write $X<Y$ if there is an
injection $X\hookrightarrow Y$ but no bijection $X\to Y$. The dual
Cantor--Schr\"oder--Bernstein principle $\mathsf{CSB}^{*}$ asserts that
$X\leq^{*}Y$ and $Y\leq^{*}X$ imply a bijection $X\to Y$; the Weak
Partition Principle $\mathsf{WPP}$ asserts that $X\leq^{*}Y$ implies
$\neg(Y<X)$. Both statements quantify over all sets. If either set in
the dual-CSB hypothesis is empty, so is the other.

In $\ZF$ one has
\[
 \PP\ \Longrightarrow\ \mathsf{CSB}^{*}\ \Longrightarrow\ \mathsf{WPP}.
\]
Indeed, the empty case of $X\leq^{*}Y\Rightarrow X\hookrightarrow Y$
has the empty injection, so the first implication follows by applying
$\PP$ twice and ordinary Cantor--Schr\"oder--Bernstein. For the second,
if $X\leq^{*}Y$ and $Y<X$, the injection $Y\hookrightarrow X$ gives
$Y\leq^{*}X$: for nonempty $Y$, invert it on its image and extend by
one fixed element of $Y$; for empty $Y$, use the definition. Then
$\mathsf{CSB}^{*}$ gives a bijection, contrary to $Y<X$. See
\cite[Definitions~2.1--2.2, Proposition~2.3, and the proofs of
Corollary~8.2 and Proposition~9.8]{BlassKulshreshtha2026}.

These $\ZF$ implications hold internally in each model. Thus
Theorem~\ref{thm:main-consistency} also gives
\[
 \Con(\ZF)\ \Longrightarrow\
 \Con(\ZF+\PP+\ACwo+\mathsf{CSB}^{*}+\mathsf{WPP}+\neg\AC).
\]
In particular, assuming $\Con(\ZF)$, neither
$\mathsf{CSB}^{*}\Rightarrow\AC$ nor $\mathsf{WPP}\Rightarrow\AC$ is
provable in $\ZF$. In the separate transitive setting of
Theorem~\ref{thm:main-transitive}, the same $W$ also satisfies
$\mathsf{CSB}^{*}$ and $\mathsf{WPP}$. The assertion $\ACwo$ remains
true there but is not needed for the implication chain.
\end{remark}

\subsection{An explicit family without a choice function}

The nonwellorderability of the label set already witnesses failure of
$\AC$.  The following elementary argument also identifies a family of
nonempty sets with no choice function.

\begin{lemma}\label{lemma:explicit-choice-failure}
In $\ZF$, if a set $A$ is not wellorderable, then
$\mathcal P(A)\setminus\{\varnothing\}$ has no choice function.
\end{lemma}

\begin{proof}
Suppose $c$ were a choice function for this family, and let $\lambda$
be the Hartogs ordinal of $A$, so there is no injection
$\lambda\longrightarrow A$.  Foundation gives $A\notin A$.
By transfinite recursion define $a:\lambda\longrightarrow A\cup\{A\}$
as follows.  At stage $\alpha$, put
$S_\alpha=A\setminus\ran(a|\alpha)$ and set
\[
 a(\alpha)=
 \begin{cases}
 c(S_\alpha),&S_\alpha\ne\varnothing,\\
 A,&S_\alpha=\varnothing.
 \end{cases}
\]
If every $S_\alpha$ were nonempty, the recursion would choose a
different member of $A$ at every stage, giving an injection
$\lambda\longrightarrow A$.  Hence there is a least $\alpha<\lambda$
with $S_\alpha=\varnothing$.  Before that stage no value is the
sentinel $A$, and $a|\alpha$ is a bijection onto $A$.  This
wellorders $A$, a contradiction.  Hartogs' theorem and the recursion
used here are available in $\ZF$.
\end{proof}

\subsection{The resulting symmetric model}

We return to the construction.  The cardinal choices in the next
statement are computed in the ground model.  They allow the same
argument to be used over a given choice ground as well as over a
constructible ground in the consistency deduction.

Over a general ground, the statement is read in its internal forcing
interpretation. Ground objects are identified with their canonical
check-name images, and ordinals are those of this interpretation.
In the sequence assertion, $\omega$ is the ground's internal
natural-number object, and $z\in V$ means that $z$ is represented
by a ground check name. When $V$ is externally transitive, these
are the usual evaluated objects and literal inclusions. The two
interpretations are distinguished in Theorem~\ref{thm:symmetric-model}
and constructed separately in Section~\ref{sec:metatheory}.

\begin{theorem}[Construction summary]\label{thm:construction-summary}
Let $V$ be a ground model of $\ZFC$, put
\[
 \kappa=\omega_1^V,\qquad
 \mu=\bigl((2^{\aleph_0})^+\bigr)^V,
 \qquad |I|^V=\mu^+,
\]
and use the monoid of Theorem~\ref{thm:monoid} and the symmetric
system constructed above.  In an ordinary generic extension for
its product forcing, the resulting symmetric model $W$ satisfies
\[
 W\models\ZF+\PP+\ACwo+\neg\AC.
\]
It contains $V$ and has the same ordinals as $V$.  It adds no
countable sequences of ground elements: if $b\in V$ and
$z\in W$ is a function $\omega\longrightarrow b$, then $z\in V$.
The set of labels $A$ is not wellorderable in $W$, and
\[
 \mathcal P^W(A)\setminus\{\varnothing\}
\]
is an explicit family in $W$ with no choice function in $W$.
\end{theorem}

\begin{proof}
The symmetric system is for one set forcing, with the set group and
normal filter specified above.  The set-forcing symmetric-model
theorem, in its internal form, therefore gives full $\ZF$, the
canonical ground embedding, and preservation of the interpreted
ordinals.  For an externally transitive ground its evaluated
hereditarily symmetric interpretation is transitive.  These are
pure-set models, so Foundation and the full Separation and Replacement
schemata are included.

Proposition~\ref{prop:presentation} gives the first hypothesis of
Lemma~\ref{lemma:pp-reduction} inside $W$, and
Proposition~\ref{prop:ordinal-choice} gives the second, for all
ordinal-indexed families in $W$.  Theorem~\ref{thm:quotient-classification},
together with Remark~\ref{rem:quotients-versus-images}, gives the third.
The lemma proves $\PP$ for every surjection in $W$, with its witnessing
injection in $W$.

Proposition~\ref{prop:not-choice} proves that $A$ is not wellorderable
in $W$.  Applying Lemma~\ref{lemma:explicit-choice-failure} inside
$W$ gives the displayed family and failure of $\AC$ there.
Finally, Proposition~\ref{prop:no-countable-sequences} supplies the
assertion about countable sequences.  This last assertion concerns
the symmetric model itself; it does not require countable closure
of the full product forcing.
\end{proof}

The statement of Theorem~\ref{thm:construction-summary} concerns the
model produced by the forcing construction.  Extracting an ordinary
relative-consistency implication, and obtaining an externally
transitive model under the stronger model-existence hypothesis,
are separate deductions given next.
